\section{Differential classes of a surface algebra}
\label{ce:surface}

A finite coefficient algebra can support a nonzero central term in a
current Lie algebra.  The simplest diagnostic uses two commuting Cartan
currents.  If $h=\operatorname{diag}(1,-1)\in\mathfrak{sl}_2$ and $u,v$
are degree-two coefficients, the central extension constructed below has
\[
 [h\otimes u,h\otimes v]=2\overline{u\,dv}.
\]
The ordinary current bracket on this pair is zero.  Thus the question is
whether $u\,dv$ survives modulo exact differentials.  For a coefficient
algebra with a $22$-dimensional degree-two component and a nonzero
quadratic pairing, these classes form a vector space of dimension $231$.
This counts central directions in an algebraic current
extension; interpreting them as operators on a space of physical states
requires a representation of that extension.

Let $\kk$ be a field of characteristic zero.  All tensor products and
differentials in this section are over $\kk$.  Let $V$ be a vector space of
finite dimension $n$, placed in internal degree $2$, and let
$q:V\times V\to \kk$ be a symmetric bilinear form.  Define the unital graded
commutative algebra
\begin{equation}\label{ce:algebra}
 A=\kk1\oplus V\oplus \kk\omega,\qquad
 |\omega|=4,\qquad uv=q(u,v)\omega,\qquad V\omega=\omega^2=0.
\end{equation}
Associativity follows because every product of three positive-degree
elements is zero.  The symbol $\omega$ is an independent nonzero basis
element, including when $q=0$ or $V=0$.

The module of K\"ahler differentials $\Omega^1_{A/\kk}$ represents
$\kk$-linear derivations into $A$-modules.  Its universal derivation
$d:A\to\Omega^1_{A/\kk}$ has internal degree zero:
$|da|=|a|$ and $d(ab)=a\,db+b\,da$.  This is a degree-preserving
universal derivation, rather than a cochain differential of degree one.
We write
\begin{equation}\label{ce:quotient}
 Z(A)=\Omega^1_{A/\kk}/dA,\qquad
 dA=\{da:a\in A\},
\end{equation}
where the quotient is taken in graded $\kk$-vector spaces.  The complete
presentation and its universal property will be proved directly below.

The case $V=0$ already separates the two carriers.  Here
$A=\kk[\omega]/(\omega^2)$ and
\[
 \Omega^1_{A/\kk}=A\,d\omega/(2\omega\,d\omega)=\kk\,d\omega\ne0,
 \qquad dA=\kk\,d\omega,\qquad Z(A)=0.
\]
Vanishing modulo exact differentials does not imply vanishing of the
differential module.

\begin{proposition}[The complete differential presentation]
\label{ce:presentation}
The $A$-module $\Omega^1_{A/\kk}$ is generated by the linear symbols $dv$,
$v\in V$, and $d\omega$, with precisely the $A$-linear relations
\begin{align}
 u\,dv+v\,du&=q(u,v)d\omega, &&u,v\in V,\label{ce:rel4}\\
 v\,d\omega+\omega\,dv&=0, &&v\in V,\label{ce:rel6}\\
 2\omega\,d\omega&=0.\label{ce:rel8}
\end{align}
This presentation is valid over every field, with the integer coefficients
interpreted in that field.
\end{proposition}

\begin{proof}
Choose a basis $v_1,\ldots,v_n$ of $V$.  The polynomial presentation of $A$
has generators $v_1,\ldots,v_n,\omega$ and defining relations
$v_iv_j-q(v_i,v_j)\omega$, $v_i\omega$, and $\omega^2$.
These relations reduce every monomial to a linear combination of
$1,v_1,\ldots,v_n,\omega$.  The map from the presented algebra to
\eqref{ce:algebra} sends these elements to a basis, so there are no further
relations.

Let $M$ be the module presented in the statement.  Assigning arbitrary
images to $dv_i,d\omega$ defines a derivation from the polynomial algebra
to an $A$-module.  It factors through $A$ exactly when it kills the
derivatives of the displayed defining relations, which are
\eqref{ce:rel4}--\eqref{ce:rel8}.  Derivatives of every element of their
ideal then vanish: $D(fr)=fD(r)+rD(f)=0$.  Thus $M$ has the defining
universal property of $\Omega^1_{A/\kk}$.  This argument also proves the
assertion in arbitrary characteristic.
\end{proof}

For a uniform description of degree six, set
\begin{equation}\label{ce:Tq}
 T_q=V\big/\left\langle
 q(u,v)w+q(v,w)u+q(w,u)v:u,v,w\in V
 \right\rangle.
\end{equation}
The class of $v$ in $T_q$ will represent $\omega\,dv$.

\begin{theorem}[All internal degrees]
\label{ce:differentials}
In characteristic zero the only possibly nonzero components of
$\Omega^1_{A/\kk}$ and $Z(A)$ are
\[
\begin{array}{c|ccc}
 j&2&4&6\\ \hline
 (\Omega^1_{A/\kk})^j&V&\bigwedge^2V\oplus \kk\,d\omega&T_q\\
 Z(A)^j&0&\bigwedge^2V&T_q.
\end{array}
\]
The vector space $T_q$ is zero if $q\ne0$, and equals $V$ if $q=0$.
In degree four the canonical quotient identification is
\begin{equation}\label{ce:wedge}
 \bigwedge^2V\xrightarrow{\ \sim\ }Z(A)^4,
 \qquad u\wedge v\longmapsto\overline{u\,dv}.
\end{equation}
In particular, nondegeneracy of $q$ is unnecessary for the conclusion
$Z(A)=Z(A)^4\cong\bigwedge^2V$; nonvanishing of $q$ suffices.
\end{theorem}

\begin{proof}
The free module before imposing relations has components only in
degrees $2,4,6,8$.  In degree two there are no relations, giving $V$.
In degree four its presentation is
\[
 (V\otimes V\oplus \kk\,d\omega)
 \big/\langle u\otimes v+v\otimes u-q(u,v)d\omega\rangle.
\]
An explicit isomorphism to $\bigwedge^2V\oplus \kk\,d\omega$ sends
$u\,dv$ to $u\wedge v+\tfrac12q(u,v)d\omega$ and fixes $d\omega$.
Its inverse sends $u\wedge v$ to
$\tfrac12(u\,dv-v\,du)$ and fixes $d\omega$.  These two formulas
compose to the identity by \eqref{ce:rel4}.

Degree six starts with the two copies $V\,d\omega$ and $\omega\,dV$.
Relation \eqref{ce:rel6} eliminates the first.  Multiplying
\eqref{ce:rel4} by $w\in V$ then gives
\[
 q(w,u)\omega\,dv+q(w,v)\omega\,du+q(u,v)\omega\,dw=0.
\]
These are precisely the relations defining $T_q$.  There are no other
degree-six relations in the module presentation.  In degree eight the
only generator is $\omega\,d\omega$, killed by \eqref{ce:rel8}.
Higher degrees have no generators.  Multiples of the relations in
degrees eight and above cannot impose additional relations in lower
degrees because the presentation is homogeneous and $A$ is nonnegatively
graded.

If $q=0$, the relations in \eqref{ce:Tq} are zero.  If $q\ne0$, symmetry
and invertibility of $2$ imply that some $e\in V$ has $q(e,e)\ne0$:
otherwise $q(u+v,u+v)-q(u,u)-q(v,v)=2q(u,v)$ would vanish for all $u,v$.
Write $t_v$ for the class of $v$ in $T_q$.  The triples $(e,e,e)$ and
$(e,e,v)$ give
\[
 3q(e,e)t_e=0,\qquad q(e,e)t_v+2q(e,v)t_e=0.
\]
Hence $t_e=0$ and then $t_v=0$ for every $v$.

Finally, $dA$ consists of $dV$ in degree two and $\kk\,d\omega$ in
degree four.  Both maps from these summands of $A$ are injective by the
explicit descriptions above.  There are no elements of $A$ in degrees
six or eight whose derivatives could contribute.  Taking the vector
space quotient proves the table and \eqref{ce:wedge}.
\end{proof}

The module action can also be read without choosing a basis.  Put
$s_{u,v}=\tfrac12(u\,dv-v\,du)$ and $t_v=\omega\,dv$.
Besides $w\cdot dv=s_{w,v}+\tfrac12q(w,v)d\omega$, its only potentially
nonzero positive-degree actions are
\[
 w\cdot d\omega=-t_w,\qquad \omega\cdot dv=t_v,\qquad
 w\cdot s_{u,v}=\tfrac12\bigl(q(w,u)t_v-q(w,v)t_u\bigr).
\]
All remaining such actions vanish by degree.  Together with
\eqref{ce:Tq} these formulas describe the full $A$-module, as well as its
graded underlying vector space.

\begin{example}[Small dimensions and a degenerate form]
\label{ce:small}
If $n=1$ and $q(v,v)=c\ne0$, then
$A\cong \kk[v]/(v^3)$ by $\omega=c^{-1}v^2$.  Its differential module is
$A\,dv/(3v^2dv)$, with basis $dv,v\,dv$ in degrees two and four.
Here $d\omega=2c^{-1}v\,dv$ and $Z(A)=0$.
If instead $q=0$ and $n=1$, then
\[
 \Omega^1_{A/\kk}=\kk\,dv\oplus \kk\,d\omega\oplus \kk\,\omega\,dv,
 \qquad Z(A)=\kk\,\overline{\omega\,dv}
\]
in degrees $2,4,6$ and $6$, respectively.  The algebra has two independent
square-zero generators $v,\omega$, so it is not $\kk[v]/(v^3)$.
For $n=2$ and Gram matrix $\operatorname{diag}(1,0)$ the form is
degenerate but nonzero: $Z(A)$ is one-dimensional in degree four and
has no degree-six component.  For the zero form in dimension two its
degree-four and degree-six dimensions are $1$ and $2$.
\end{example}

\begin{remark}[Why the quotient is a vector space]
If $n\ge2$, take independent $u,v$.  Then $dv\in dA$, whereas
$u\,dv\notin dA$ by \eqref{ce:wedge}.  Thus multiplication by $u$
does not descend to \eqref{ce:quotient}.  When $n=1$ and $q=0$,
$d\omega\in dA$ but $v\,d\omega=-\omega\,dv\notin dA$ gives the same
failure.  In general $dA$ generates $\Omega^1_{A/\kk}$ as an $A$-module;
therefore $dA$ is an $A$-submodule if and only if $Z(A)=0$.  An
$A$-module quotient by the submodule generated by exact differentials
would always be zero and would lose every central class under discussion.
\end{remark}

\begin{remark}[The role of the characteristic]
The calculation of the differential module uses only invertibility of
$2$ and $3$.  Theorem~\ref{ce:differentials} therefore holds over any
field of characteristic different from $2,3$.  In characteristic $3$,
$\kk[v]/(v^3)$ with $|v|=2$ has no differential relation $v^2dv=0$;
its quotient by exact differentials contains the nonzero degree-six
class $\overline{v^2dv}$.  In characteristic $2$ the example
$\kk[\omega]/(\omega^2)$ has the additional nonzero degree-eight class
$\overline{\omega\,d\omega}$.  Thus both excluded primes change the
answer.  Proposition~\ref{ce:presentation}, with no division, still
specifies all relations in these characteristics.  Characteristic zero
will also be used below for polynomial integration and for the
primitive calculation in an enveloping algebra.
\end{remark}

\section{The current cocycle and its central directions}
\label{ce:current}

For the sign calculation, allow $B$ to be any unital
$\mathbb Z$-graded supercommutative $\kk$-algebra.  The parity of a
homogeneous element is its internal degree modulo two.  Define
$\Omega^1_{B/\kk}$ using an even universal derivation: take the free
graded left $B$-module on symbols $db$ of degree $|b|$, impose
$\kk$-linearity, $d1=0$, and the Leibniz relations displayed below.
A degree-zero module map from this quotient is exactly an even
derivation into the target module, which proves its universal property.
For a homogeneous
module element $m$, its right action is $mb=(-1)^{|m||b|}bm$, so
\[
 d(ab)=(da)b+a\,db=(-1)^{|a||b|}b\,da+a\,db.
\]
The vector space $Z(B)=\Omega^1_{B/\kk}/dB$ retains the internal grading.
Let $\mathfrak g$ be an ordinary Lie algebra, placed in degree zero,
with a symmetric invariant bilinear form $\kappa$:
\begin{equation}\label{ce:kappa}
 \kappa([x,y],z)=\kappa(x,[y,z]),\qquad
 \kappa(x,y)=\kappa(y,x).
\end{equation}
Thus all super signs here come from $B$; no element of $\mathfrak g$ has
odd parity.  No finiteness or nondegeneracy assumption on $\kappa$ is
needed for the next construction.

\begin{theorem}[Central current extension]
\label{ce:cocycle}
The graded vector space
\[
 \widehat L_B=(\mathfrak g\otimes B)\oplus Z(B)
\]
is a Lie superalgebra with $Z(B)$ central and, for homogeneous $a,b$,
\begin{equation}\label{ce:bracket}
 [x\otimes a,y\otimes b]
 =[x,y]\otimes ab+\kappa(x,y)\overline{a\,db}.
\end{equation}
Both terms have internal degree $|a|+|b|$.  Projection gives an exact
sequence of Lie superalgebras
\[
 0\longrightarrow Z(B)\longrightarrow\widehat L_B
 \longrightarrow\mathfrak g\otimes B\longrightarrow0.
\]
\end{theorem}

Here the hat in $\widehat L_B$ denotes this central extension.
Its carrier is the displayed algebraic direct sum; no height completion
is being taken.

\begin{proof}
Exactness of $d(ab)$ gives
$\overline{a\,db}=-(-1)^{|a||b|}\overline{b\,da}$, which proves graded
skew symmetry using \eqref{ce:kappa} and supercommutativity of $B$.
Write $p=|a|$, $q'=|b|$, $r=|c|$ modulo two.  In the signed cyclic
Jacobi sum
\[
 (-1)^{pr}[[x\otimes a,y\otimes b],z\otimes c]
       +\text{the two cyclic terms},
\]
the noncentral part vanishes by the ordinary Jacobi identity in
$\mathfrak g$: rewriting every coefficient in the order $abc$ gives
the same sign $(-1)^{pr}$ on its three terms.

Symmetry and invariance give a common scalar
$f=\kappa([x,y],z)=\kappa([y,z],x)=\kappa([z,x],y)$ in the central
part.  That part divided formally into its scalar and vector factors is
$f$ times
\begin{equation}\label{ce:jacobi}
 (-1)^{pr}\overline{ab\,dc}
 +(-1)^{pq'}\overline{bc\,da}
 +(-1)^{q'r}\overline{ca\,db}
 =(-1)^{pr}\overline{d(abc)}=0.
\end{equation}
To check the displayed equality, expand
$d(abc)=(da)bc+a(db)c+ab\,dc$, and move $da$ and $db$ to the right.
The first two coefficients acquire signs $(-1)^{p(q'+r)}$ and
$(-1)^{pr+q'r}$ when written as $bc\,da$ and $ca\,db$.
This proves the required identity without assuming $f\ne0$.
Every Jacobi identity involving a central argument vanishes separately.
\end{proof}

For the algebra $A$ in \eqref{ce:algebra} all parities are even, so
\eqref{ce:bracket} is an ordinary Lie bracket with an additional internal
grading.  If $q\ne0$, its central contribution is nonzero only when
$a,b$ both have degree two.  Under \eqref{ce:wedge} it is simply
$\kappa(x,y)(a\wedge b)$.

\section{Universal extensions of ordinary currents}
\label{ce:universal-section}

The cocycle theorem constructs an extension.  Universality asks whether
it accounts for every possible central term.  For this question, take
$B$ to be an ordinary commutative unital $\kk$-algebra and specialize to
$\mathfrak g=\mathfrak{sl}_2(\kk)$, with the trace pairing.  The field
$\kk$ still has characteristic zero.  All tensor products are algebraic,
with no topology or completion.  An ungraded $B$ is regarded as
concentrated in degree zero.  The universal result in this section
concerns ordinary Lie algebras; it does not extend the preceding
super-coefficient construction by assumption.

Fix the trace-zero matrices
\[
 e=\begin{pmatrix}0&1\\0&0\end{pmatrix},\qquad
 f=\begin{pmatrix}0&0\\1&0\end{pmatrix},\qquad
 h=\begin{pmatrix}1&0\\0&-1\end{pmatrix}.
\]
The commutator and trace give
\begin{equation}\label{ce:uc:sl2}
 [h,e]=2e,\qquad[h,f]=-2f,\qquad[e,f]=h,
\end{equation}
and $\kappa(e,f)=1$, $\kappa(h,h)=2$; the pairings of $e$ with itself,
$f$ with itself, and $h$ with $e$ or $f$ vanish.  The current Lie
algebra is $L_B=\mathfrak{sl}_2(\kk)\otimes B$, with
$[x\otimes a,y\otimes b]=[x,y]\otimes ab$.  Write
$e_a=e\otimes a$, and similarly $f_a,h_a$.

A central extension is a surjective Lie map $p:E\to L_B$ whose
kernel $W$ is central, meaning $[W,E]=0$.  It is universal if every
central extension of $L_B$ receives a unique Lie map from it over
$L_B$.  No finite-dimensionality assumption is placed on $B$ or $W$.

Write $\langle a,b\rangle=\overline{a\,db}\in Z(B)$.
The differential presentation above gives the complete vector-space
relations
\begin{equation}\label{ce:uc:symbol-relations}
 \langle1,b\rangle=0,\qquad
 \langle a,bc\rangle=\langle ab,c\rangle+\langle ac,b\rangle,
\end{equation}
with bilinearity in $a,b$.  The second relation is every $B$-multiple
of the product rule; the first kills exactly $dB$.  Taking $b=c=1$
also gives $\langle a,1\rangle=0$, so no relation for $d1$ is missing.
Taking $a=1$ gives $\langle a,b\rangle=-\langle b,a\rangle$.
Thus the extension of Theorem~\ref{ce:cocycle} is
$\widehat L_B=L_B\oplus Z(B)$ with projection $\widehat p$, and
\begin{equation}\label{ce:uc:cartan-witness}
 [(h_a,0),(h_b,0)]=(0,2\langle a,b\rangle).
\end{equation}
Every central differential class is detected by a pair of currents
whose bracket in $L_B$ vanishes.  The classification below shows that
these pairs detect all central cocycles up to changes of lifts.

\subsection{Changing lifts and normalizing every cocycle}
\label{ce:uc:normalization}

Choose a linear section $s:L_B\to E$ of a central extension $p$.
Such a section exists by lifting a vector-space basis.  Identify $W$
with its image in $E$.  The failure of $s$ to preserve brackets is
\[
 c(X,Y)=[s(X),s(Y)]-s([X,Y])\in W.
\]
It is alternating and bilinear.  Jacobi in $E$, centrality of $W$, and
Jacobi in $L_B$ imply
\begin{equation}\label{ce:uc:cocycle}
 c([X,Y],Z)+c([Y,Z],X)+c([Z,X],Y)=0.
\end{equation}
An alternating bilinear map satisfying this equation is called a
\emph{central two-cocycle}.  Conversely any such map defines an extension
on $L_B\oplus W$ by $[(X,u),(Y,v)]=([X,Y],c(X,Y))$; the same expansion
proves its Jacobi identity.

Replacing $s$ by $s+\beta$, for a linear map $\beta:L_B\to W$, replaces
$c$ by $c-\beta\circ[-,-]$.  Maps $\beta\circ[-,-]$ are called
\emph{coboundaries}; they satisfy \eqref{ce:uc:cocycle} because the Lie
bracket satisfies Jacobi.  Define $H^2_{\mathrm{Lie}}(L_B,W)$ to be the vector space
of these cocycles modulo coboundaries.  This definition concerns a
trivial coefficient action: elements of $L_B$ act as zero on $W$.
No further Lie-cohomology results are needed.  The same cocycle
definition and section-change calculation apply to any ordinary Lie
algebra in place of $L_B$.

\begin{theorem}[Explicit classification of all central cocycles]
\label{ce:uc:classification}
For every vector space $W$ and every central two-cocycle
$c:L_B\times L_B\to W$, there are unique linear maps
$\beta:L_B\to W$ and $\lambda:Z(B)\to W$ such that
\begin{equation}\label{ce:uc:classification-formula}
 c(x\otimes a,y\otimes b)
 =\beta([x,y]\otimes ab)+\kappa(x,y)\lambda(\langle a,b\rangle).
\end{equation}
They are given explicitly by
\begin{align}
 \beta(e_a)&=\tfrac12c(h_1,e_a),&
 \beta(f_a)&=-\tfrac12c(h_1,f_a),&
 \beta(h_a)&=c(e_1,f_a),\label{ce:uc:beta}\\
 \lambda(\langle a,b\rangle)&=\tfrac12c(h_a,h_b).
 &&&\label{ce:uc:lambda}
\end{align}
Consequently the map
$\operatorname{Hom}_\kk(Z(B),W)\to H^2_{\mathrm{Lie}}(L_B,W)$ induced by the central
term in \eqref{ce:bracket} is an isomorphism.
\end{theorem}

\begin{proof}
First insert $(h_1,e_1,f_a)$ in \eqref{ce:uc:cocycle}.  The first and third
terms are $2c(e_1,f_a)$ and $2c(f_a,e_1)$, which cancel.  The remaining
term is $c(h_a,h_1)$, so
\begin{equation}\label{ce:uc:h1h}
 c(h_1,h_a)=0\qquad(a\in B).
\end{equation}
Define $\beta$ by \eqref{ce:uc:beta}, using
$L_B=e_B\oplus f_B\oplus h_B$ as a vector space, and put
$c_0=c-\beta\circ[-,-]$.  The relations \eqref{ce:uc:sl2} and
\eqref{ce:uc:h1h} give
\begin{equation}\label{ce:uc:normalized}
 c_0(h_1,X)=0\quad(X\in L_B),\qquad c_0(e_1,f_a)=0.
\end{equation}
In particular these normalizations impose no extra hypothesis on $c$.

The adjoint action of $h_1$ has eigenvalues $2,-2,0$ on $e_B,f_B,h_B$.
If $[h_1,X]=rX$ and $[h_1,Y]=sY$, the cocycle equation on $(h_1,X,Y)$
and \eqref{ce:uc:normalized} give
\[
 (r+s)c_0(X,Y)=0.
\]
The nonzero sums are $2,-2,4,-4$, all invertible in $\kk$.
Therefore the only potentially nonzero entries are
\[
 P(a,b)=c_0(e_a,f_b),\qquad R(a,b)=c_0(h_a,h_b).
\]
Here $R$ is alternating and $P(1,b)=0$.

The cocycle equation on $(e_a,f_b,h_c)$ reads
\begin{equation}\label{ce:uc:decisive}
 R(ab,c)-2P(a,bc)+2P(ac,b)=0.
\end{equation}
Setting $a=1$ gives $R(b,c)=-2P(c,b)$.  Alternation of $R$ implies
$P(c,b)=-P(b,c)$ and hence $R(b,c)=2P(b,c)$.
Substitution back into \eqref{ce:uc:decisive} yields
\begin{equation}\label{ce:uc:P-relations}
 P(1,b)=0,\qquad P(a,bc)=P(ab,c)+P(ac,b).
\end{equation}
These are precisely the complete differential-symbol relations
\eqref{ce:uc:symbol-relations}.  There is therefore a unique linear map
$\lambda:Z(B)\to W$ with $\lambda(\langle a,b\rangle)=P(a,b)$.
It also satisfies \eqref{ce:uc:lambda}, because the coboundary correction
does not change $c(h_a,h_b)$: the current bracket of these elements is
zero.

The listed values of the trace pairing show that $c_0$ equals
$\kappa(x,y)\lambda(\langle a,b\rangle)$ on every pair of the three
basis types.  Bilinearity proves \eqref{ce:uc:classification-formula}.
Conversely this formula defines a cocycle for arbitrary $\beta,\lambda$,
by the coboundary identity and the proved bracket
\eqref{ce:bracket}.

For uniqueness, any expression \eqref{ce:uc:classification-formula}
evaluated on $(h_a,h_b)$ forces \eqref{ce:uc:lambda}, since $2$ is
invertible and the symbols span $Z(B)$.  Evaluation on
$(h_1,e_a),(h_1,f_a),(e_1,f_a)$ then forces \eqref{ce:uc:beta}, because
$\langle1,a\rangle=0$.  Thus the decomposition is unique.
Quotienting this decomposition by coboundaries proves the final
isomorphism, including for infinite-dimensional $B$ or $W$.
\end{proof}

The scalar case $W=\kk$ gives more than a lower bound on the number of
central classes.  It proves that no additional classes exist.  The
central error can be read off from commuting Cartan currents by
\eqref{ce:uc:lambda}, and the three corrections in \eqref{ce:uc:beta} remove
everything else.

\subsection{The universal map and its uniqueness}
\label{ce:uc:universal-section}

\begin{theorem}[Universal central extension]
\label{ce:uc:universal}
For every commutative unital $\kk$-algebra $B$, the projection
$\widehat p:\widehat L_B\to L_B$ of \eqref{ce:bracket} is
a universal central extension in the category of ordinary Lie algebras
over $\kk$.
\end{theorem}

\begin{proof}
Let $p:E\to L_B$ be any central extension with kernel $W$, and choose
a linear section $s$.  Apply Theorem~\ref{ce:uc:classification} to its
cocycle, obtaining $\beta$ and $\lambda$.  Set $s_0=s+\beta$, with
$\beta$ regarded as a map into the central subspace $W\subset E$.
The calculation of section change gives
\[
 [s_0(x\otimes a),s_0(y\otimes b)]
 -s_0([x,y]\otimes ab)=\kappa(x,y)\lambda(\langle a,b\rangle).
\]
It follows that the linear map
\begin{equation}\label{ce:uc:factor}
 F:\widehat L_B\longrightarrow E,\qquad
 F(X,z)=s_0(X)+\lambda(z)
\end{equation}
preserves brackets and satisfies $pF=\widehat p$.  This constructs the
required factorization for every central extension, with no dimension
assumption on its kernel.

For uniqueness, a Lie algebra is called \emph{perfect} if it is spanned
by its brackets.  The extension $\widehat L_B$ is perfect.  Indeed
\[
 [(h_1,0),(e_a,0)]=(2e_a,0),\qquad
 [(h_1,0),(f_a,0)]=(-2f_a,0),\qquad
 [(e_1,0),(f_a,0)]=(h_a,0),
\]
because $\langle1,a\rangle=0$.  These span its current summand.
Equation~\eqref{ce:uc:cartan-witness} spans its central summand as well.

If $F,G:\widehat L_B\to E$ are Lie maps over $L_B$, then
$(F-G)(X)\in W$ for every $X$.  Since $W$ is central,
\[
 (F-G)([X,Y])=[F(X),F(Y)]-[G(X),G(Y)]=0.
\]
Perfectness forces $F=G$.  In particular the map in \eqref{ce:uc:factor}
is independent of the initially chosen section.
\end{proof}

The construction can be expressed without leaving the vector space
$L_B\oplus W$.  For any $\lambda:Z(B)\to W$, give
$E_\lambda=L_B\oplus W$ the bracket
\[
 [(X,u),(Y,v)]=([X,Y],\lambda\sigma(X,Y)),\qquad
 \sigma(x\otimes a,y\otimes b)=\kappa(x,y)\langle a,b\rangle.
\]
This is the extension obtained by replacing each universal central
symbol with its image under $\lambda$.  The map
$\widehat L_B\to E_\lambda$, $(X,z)\mapsto(X,\lambda z)$, need not
be onto: $W$ may contain central directions unused by the bracket.
For a general $E$, the map $(X,w)\mapsto s_0(X)+w$ is an isomorphism
$E_\lambda\to E$, compatible with both the kernel and quotient.
Thus the classification describes all extensions, including those with
unused central summands, while universality gives a unique map to each.

\begin{proposition}[Functoriality and grading]
\label{ce:uc:functor}
A unital algebra map $f:B\to B'$ induces Lie maps
\[
 L_f(x\otimes a)=x\otimes f(a),\qquad
 \widehat L_f(x\otimes a,z)
  =(x\otimes f(a),Z(f)z),
\]
where $Z(f)\langle a,b\rangle=\langle f(a),f(b)\rangle$.
These maps preserve identities and compositions.  The classification
isomorphism is covariant in $W$ and contravariant in $B$, by composition
with a map $W\to W'$ or pullback along $L_f$, respectively.

If $B$ is graded with even degrees, give $db$ degree $|b|$ and
$\langle a,b\rangle$ degree $|a|+|b|$.  The same extension is universal
among graded central extensions with degree-zero maps.
\end{proposition}

\begin{proof}
The algebra map preserves both relations \eqref{ce:uc:symbol-relations},
so $Z(f)$ is well-defined.  Applying $f$ to products and to symbols in
\eqref{ce:bracket} proves the Lie assertion.  The formulas on
symbols prove identity and composition laws.  The cocycle formulas
commute with maps of $W$ and with substitution of $f(a)$ for $a$,
proving the two variance assertions.

For the graded statement choose lifts degree by degree to obtain a
homogeneous section.  The unit has degree zero, so each formula
\eqref{ce:uc:beta} preserves the degree of its argument.
Formula~\eqref{ce:uc:lambda} sends a symbol of degree $|a|+|b|$ into
that degree of $W$.  Hence \eqref{ce:uc:factor} preserves degrees.
Uniqueness among all Lie maps already implies uniqueness among
homogeneous ones.  No cochain shift or parity reversal is introduced.
\end{proof}

\subsection{Central parameters for the quadratic algebra}

Apply the universal theorem to the coefficient algebra
\eqref{ce:algebra}.  Theorem~\ref{ce:differentials} supplies its full
differential quotient, so for every vector space $W$ the classification
becomes
\begin{equation}\label{ce:all-central-parameters}
 H^2_{\mathrm{Lie}}(\mathfrak{sl}_2\otimes A,W)
 \cong\begin{cases}
 \operatorname{Hom}_\kk(\bigwedge^2V,W),&q\ne0,\\
 \operatorname{Hom}_\kk(\bigwedge^2V\oplus V,W),&q=0.
 \end{cases}
\end{equation}
The coefficient degrees of these summands are four and six.
A scalar-valued cochain supported on arguments of total degree $j$
has internal degree $-j$.  An arbitrary sum of the two scalar sectors
need not be homogeneous.

More explicitly, write $a=a_01+u+s\omega$ and $b=b_01+v+t\omega$.
An alternating map $\Lambda:V\times V\to W$, together with a linear
map $T:V\to W$ when $q=0$, gives
\begin{equation}\label{ce:uc:all-parameters}
 \lambda(\langle a,b\rangle)=
 \begin{cases}
 \Lambda(u,v),&q\ne0,\\
 \Lambda(u,v)+sT(v)-tT(u),&q=0.
 \end{cases}
\end{equation}
Terms with a unit are exact, and
$\langle u,\omega\rangle=-\langle\omega,u\rangle$, which proves
the formula.  Together with the change of lifts
\eqref{ce:uc:beta}, this describes every central cocycle.
For $V=0$, the module $\Omega^1_{A/\kk}=\kk\,d\omega$ is nonzero but
$Z(A)=0$, so the current algebra itself is its universal central
extension.  The same holds for $\dim V=1$ and $q\ne0$.
For the zero form in dimension one, the additional parameter is the
degree-six class $\langle\omega,v\rangle$.  Thus degeneracy changes
the central parameter space, not the existence of the universal map.

The classification above concerns $\mathfrak{sl}_2$ with its normalized
trace pairing.  For other ordinary Lie algebras a symmetric invariant
pairing still detects nontrivial central directions under the following
condition, without claiming their exhaustiveness.

\begin{proposition}[Every differential direction is essential]
\label{ce:essential}
Let $B$ be an ordinary commutative unital $\kk$-algebra and let
$\mathfrak g$ be an ordinary Lie algebra with a symmetric invariant
pairing $\kappa$.  Assume $\kappa(x_0,x_0)\ne0$ for some $x_0\in\mathfrak g$.  Then the map
\[
 Z(B)^*\longrightarrow H^2_{\mathrm{Lie}}(\mathfrak g\otimes B,\kk),
 \quad
 \lambda\longmapsto
 \left[(x\otimes a,y\otimes b)\longmapsto
 \kappa(x,y)\lambda(\overline{a\,db})\right]
\]
is injective.  Here $\kk$ is the trivial module, the star is the algebraic
linear dual, and Lie cohomology is ordinary Chevalley--Eilenberg
cohomology, without an internal-degree restriction on cochains.
Consequently the extension cannot split when $Z(B)\ne0$, and its pushout
along any nonzero linear map $Z(B)\to W$ cannot split either.
\end{proposition}

\begin{proof}
Theorem~\ref{ce:cocycle} says that every displayed cochain is a cocycle.
The elements $x_0\otimes a$ and $x_0\otimes b$ commute in
$\mathfrak g\otimes B$.  A scalar coboundary has the form
$\delta\ell(X,Y)=-\ell([X,Y])$ and vanishes on this pair.  The displayed
cocycle takes the value
$\kappa(x_0,x_0)\lambda(\overline{a\,db})$ on the same pair.
Classes $\overline{a\,db}$ span $Z(B)$ by the definition of the
differential module.  Thus a coboundary forces $\lambda=0$.
For a nonzero map to $W$, choose a linear functional on $W$ that is
nonzero on its image.  A splitting would give a splitting after this
scalar pushout, contradicting what was just proved.
\end{proof}

\begin{corollary}[A coefficient sector of dimension $231$]
\label{ce:k3}
In \eqref{ce:algebra}, take $\kk=\mathbb Q$, $\dim V=22$, and $q\ne0$.
Then $Z(A)=\bigwedge^2V$ in internal degree four, with dimension $231$.
For $\mathfrak g=\mathfrak{sl}_2(\mathbb Q)$ and
$\kappa(x,y)=\operatorname{tr}(xy)$, this is exactly the center of
$\widehat L_A$.  The dual of this center is isomorphic to
$H^2_{\mathrm{Lie}}(\mathfrak g\otimes A,\mathbb Q)$ by scalar pushout.
The dimensions of $\widehat L_A$ in internal degrees $0,2,4$ are
$3,66,234$, respectively; all other components vanish.
\end{corollary}

\begin{proof}
Theorem~\ref{ce:differentials} gives
$\dim Z(A)=\binom{22}{2}=231$.

The trace form is symmetric and invariant, since cyclicity of trace
gives $\operatorname{tr}([x,y]z)=\operatorname{tr}(x[y,z])$.
It has $\kappa(h,h)=2$ for $h=\operatorname{diag}(1,-1)$, so
Proposition~\ref{ce:essential} applies, and
Theorem~\ref{ce:uc:classification} proves that these are all scalar
central classes.  A central element of
$\mathfrak{sl}_2\otimes A$ must commute with $\mathfrak{sl}_2\otimes1$.
Writing its coefficients in a basis of $A$ shows that each coefficient
in $\mathfrak{sl}_2$ is central, hence zero.  Therefore the current
algebra has zero center, and the center of its extension is precisely
the kernel $Z(A)$.  Counting its direct-sum components gives the stated
dimensions.
\end{proof}

For a basis $v_1,\ldots,v_{22}$ of $V$, the elements
$c_{ij}=\overline{v_i\,dv_j}$, $i<j$, form a basis of this center.
The identity
\[
 [h\otimes v_i,h\otimes v_j]=2c_{ij}
\]
expresses each central basis element as one half of a commutator.  The
center is even, has
internal degree four, and has root weight zero for the adjoint Cartan
action.  It is not a copy of the one-dimensional component $A^4=\kk\omega$:
the two groups occupy the same internal degree but have different
definitions and dimensions.  Suppose additionally that $V_{\mathbb C}$
is supplied with a decomposition into types $(2,0),(1,1),(0,2)$ of
dimensions $1,20,1$, compatible with the coefficient multiplication.
The induced decomposition of its exterior square then has types
$(3,1),(2,2),(1,3)$ with dimensions $20,191,20$: the middle dimension
is $\binom{20}{2}+1$.
No cohomological shift or parity reversal enters these counts.  With a
scalar target placed in degree zero, the resulting scalar cocycles have
internal degree $-4$.  A pushout that preserves the grading instead
places its one-dimensional central target in degree four.

\begin{example}[Conditional surface interpretation]
Let $S$ be a compact complex K3 surface for which the following rational
cohomology presentation is supplied as a hypothesis:
\[
 H^*(S,\mathbb Q)=\mathbb Q1\oplus H^2(S,\mathbb Q)\oplus\mathbb Q\omega,
 \qquad \dim H^2(S,\mathbb Q)=22,\qquad \int_S\omega=1,
\]
with all other degrees zero, and with a nondegenerate symmetric cup
pairing $q(u,v)=\int_Su\smile v$.  Use the complex orientation and no
shifts or Tate twists.  The cohomology multiplication then gives exactly
\eqref{ce:algebra}: $u\smile v=q(u,v)\omega$, and the other
positive-degree products vanish by degree.  Thus the preceding results
apply to $V=H^2(S,\mathbb Q)$ and identify the differential quotient
with $\bigwedge^2H^2(S,\mathbb Q)$ in degree four.  If the stated Hodge
decomposition of dimensions $1,20,1$ is supplied as well, its exterior
square has the dimensions just computed.  These conclusions follow from
the specified presentations; deriving those presentations from the
geometric definition of $S$ is an additional hypothesis-verification
problem and is not part of the algebraic construction.
\end{example}

\begin{remark}[Why perfectness is necessary]
The universal theorem fixes both the ordinary current algebra and its
trace pairing.  An arbitrary invariant pairing does not by itself
imply universality.  A necessary condition is perfectness.  In fact, if $E\to L$ were
universal, the two morphisms
$E\to L\oplus E/[E,E]$ given by $e\mapsto(p(e),0)$ and
$e\mapsto(p(e),\overline e)$ would have to agree.  Hence $E=[E,E]$,
and then $L=[L,L]$.  The cocycle theorem also applies to nonzero abelian
$\mathfrak g$, where $L$ is not perfect, so it cannot imply universality
as stated.

\end{remark}

\section{Positive currents and the primitive comparison}
\label{ce:positive}

The central directions above are visible to Cartan currents.  A positive
root sector behaves differently.  Suppose a Lie algebra over a
characteristic-zero field has a split root decomposition with Cartan
subalgebra $\mathfrak t$ and a chosen positive root system, and let
$\kappa$ be invariant.  For $x\in\mathfrak g_\alpha$ and
$y\in\mathfrak g_\beta$, invariance gives
\[
 0=\kappa([H,x],y)+\kappa(x,[H,y])
   =(\alpha+\beta)(H)\kappa(x,y),\qquad H\in\mathfrak t.
\]
If $\alpha+\beta\ne0$, some $H$ makes the scalar nonzero.
Positive roots never sum to zero.  Thus
$\kappa(\mathfrak n_+,\mathfrak n_+)=0$, and the cocycle restricts to
zero on $\mathfrak n_+\otimes A$.  The pullback central extension there
has the explicit Lie section $x\otimes a\mapsto(x\otimes a,0)$.

The enveloping algebra and its primitive coproduct are those of
Section~\ref{sp:section}. Theorem~\ref{sp:pbw}, with every basis vector
even, gives the ordinary ordered basis used below, and
Theorem~\ref{sp:primitives} identifies its primitive Lie algebra.

Return to the polynomial current algebra of
Remark~\ref{current-quotient-example}, now with
$\mathfrak n_+=\kk e\subset\mathfrak{sl}_2$.
Give $t$ internal degree two, so
\[
 \mathfrak a=\kk e\otimes \kk[t],\qquad
 C=U(\mathfrak a)=\kk[e_0,e_1,e_2,\ldots],\qquad e_r=e\otimes t^r.
\]
Every generator has charge one and internal degree $2r$; the Lie bracket
is zero.  The ordinary enveloping Hopf structure makes every $e_r$
primitive.  These two gradings serve different purposes: charge counts
positive-root factors, while internal degree records the coefficient.

\begin{proposition}[Coefficient specialization of positive currents]
\label{ce:specialization}
Use $\mathfrak g=\mathfrak{sl}_2(\kk)$ with the trace pairing.
Fix $v\in V$ and define the graded algebra map
$f_v:\kk[t]\to A$ by $t\mapsto v$.  There is a Hopf algebra map
\begin{equation}\label{ce:map}
 F_v:C\longrightarrow U(\widehat L_A),\qquad
 e_r\longmapsto(e\otimes v^r,0).
\end{equation}
It preserves internal degree.  Its image is a connected algebra graded
by nonnegative charge, with charge-one primitive space
$\kk e\otimes \kk[v]$.  Its kernel is
\[
 \ker F_v=
 \begin{cases}
  (e_r:r\ge1),&v=0,\\
  (e_r:r\ge2),&v\ne0,\ q(v,v)=0,\\
  (e_r:r\ge3),&q(v,v)\ne0.
 \end{cases}
\]
The image has zero intersection with the central primitive subspace
$Z(A)\subset U(\widehat L_A)$.
\end{proposition}

\begin{proof}
The formula $f_v(t^r)=v^r$ defines an algebra map because $A$ is
commutative.  The image vectors in \eqref{ce:map} commute:
$[e,e]=0$ and $\kappa(e,e)=\operatorname{tr}(e^2)=0$.
The resulting Lie map extends to an enveloping algebra map and is a
Hopf map because it sends the primitive generators to primitives.

One has $v^2=q(v,v)\omega$ and $v^3=0$.  The nonzero vectors among
$1,v,v^2$ are linearly independent, since they have distinct internal
degrees.  Thus $\kk e\otimes \kk[v]$ has dimension one, two, or three in
the respective cases.  The ordered basis of Theorem~\ref{sp:pbw},
with a basis of this abelian Lie subalgebra extended
to a basis of $\widehat L_A$, shows that its enveloping algebra embeds
and is polynomial on these independent vectors.  This proves the exact
kernel formulas: the indicated higher modes vanish, and there are no
polynomial relations among the surviving modes.  The same basis shows
that the polynomial subalgebra meets the degree-one linear space
$\widehat L_A$ precisely in $\kk e\otimes \kk[v]$; this space is disjoint
from the direct summand $Z(A)$.
\end{proof}

The restriction to one coefficient also explains why no differential
center enters this comparison.  Functoriality of universal differentials
gives a well-defined linear map
\[
 Z(\kk[t])\longrightarrow Z(A),\qquad
 \overline{a\,db}\longmapsto
 \overline{f_v(a)\,d f_v(b)}.
\]
It sends exact elements to exact elements and therefore descends to the
stated vector space quotients.  In characteristic zero,
$\Omega^1_{\kk[t]/\kk}=\kk[t]dt=d(\kk[t])$, since
$t^r dt=d(t^{r+1}/(r+1))$.  Consequently its source is zero.
Equivalently, $Z(\kk[v])=0$: the three possibilities for $\kk[v]$ are
$\kk$, $\kk[t]/(t^2)$ and $\kk[t]/(t^3)$, whose differential monomials are
exact in characteristic zero.  Antisymmetric degree-four classes require
two independent degree-two coefficients.

By Theorem~\ref{sp:primitives}, $P(C)=\mathfrak a$ and
$P(U(\widehat L_A))=\widehat L_A$. Thus $Z(A)$ is a subspace of
central primitives in the latter enveloping algebra, while
Proposition~\ref{ce:specialization} proves that none of its nonzero
vectors lies in the positive-current image.

A geometric positive Hall algebra requires its own comparison.  Even
an identification of a positive Bogomolny--Prasad--Sommerfield algebra with $U(\mathfrak n_+)$ as
an associative algebra does not supply a Hopf structure or identify its
primitive Lie algebra.  To realize the current extension one would need
a specified Lie map from $\widehat L_A$ to the primitive Lie algebra of
a chosen Hall bialgebra, preserving the bracket
\eqref{ce:bracket}, internal degrees, and the claimed central inclusion.
The positive map \eqref{ce:map} supplies only the indicated abelian
subalgebra.  Its cocycle is zero, so it cannot establish such a central
inclusion.

There is also a grading obstruction to placing the full extension in a
connected positive-charge bialgebra.  In the $\mathfrak{sl}_2$ root
grading, $e,f,h$ have charges $1,-1,0$ and $Z(A)$ has charge zero.
The relations $[h,e]=2e$ and $[e,f]=h$ preserve this grading.  In a
connected bialgebra whose charge-zero component is $\kk1$, every
charge-zero primitive is zero.  A root-charge-preserving map from the
full extension to its primitives therefore kills $h\otimes1$; the first
relation then kills $e\otimes a$, and likewise $f\otimes a$.
Thus a realization containing nonzero Cartan and central primitives
requires a larger charge-zero component and a construction beyond the
connected positive sector.

\subsection{Central actions on state spaces}

Take again $L_B=\mathfrak{sl}_2(\kk)\otimes B$ with the trace pairing,
where $B$ is ordinary commutative and unital over the characteristic-zero
field $\kk$.  All actions below are algebraic.

There is a similarly precise distinction between a central algebra and
its representation on states.  A projective current action on a
nonzero vector space $M$ is a linear assignment
$T:L_B\to\operatorname{End}_\kk(M)$ satisfying
\[
 [T(X),T(Y)]-T([X,Y])=c(X,Y)\operatorname{id}_M
\]
for scalars $c(X,Y)$.  Jacobi of operator commutators forces
\eqref{ce:uc:cocycle}; changing $T$ by $\beta(X)\operatorname{id}_M$
changes $c$ by the coboundary above.  The classification then makes
the normalized action an actual action of $\widehat L_B$, with $z$
acting by $\lambda(z)\operatorname{id}_M$.
Conversely such an action gives the displayed projective action by
restriction to $(X,0)$.  This proves the dictionary at the operator
level; it does not construct a state space for a prescribed $\lambda$.

If $M$ is finite-dimensional and nonzero, such a scalar central action
must have $\lambda=0$.  Indeed \eqref{ce:uc:cartan-witness} implies
\[
 0=\operatorname{tr}[T(h_a),T(h_b)]
   =2\dim(M)\,\lambda(\langle a,b\rangle),
\]
and the symbols span $Z(B)$.  Thus a nonzero scalar central charge
requires an infinite-dimensional state space in characteristic zero.
